% !TEX program = xelatex
%==============================================================================
%  Phân công code vấn đáp đồ án X509Certificate
%  Phần việc chi tiết: NGUYỄN TRỌNG PHÚC
%  — Lõi mật mã, Root CA & Phát hành chứng chỉ —
%
%  Biên dịch:  xelatex PhanCong_NguyenTrongPhuc.tex  (chạy 2 lần để có mục lục)
%==============================================================================
\documentclass[12pt,a4paper]{article}

% ---- Font (XeLaTeX, hỗ trợ Unicode tiếng Việt đầy đủ) -----------------------
\usepackage{fontspec}
\setmainfont{Latin Modern Roman}
\setsansfont{Latin Modern Sans}
\setmonofont{Latin Modern Mono}

% ---- Bố cục trang -----------------------------------------------------------
\usepackage[a4paper,margin=2.4cm]{geometry}
\usepackage{setspace}\setstretch{1.18}
\usepackage{parskip}
\usepackage{amssymb} % \square cho checklist
\usepackage{csquotes} % \enquote{...}

% ---- Màu sắc ----------------------------------------------------------------
\usepackage{xcolor}
\definecolor{accent}{HTML}{1F4E79}
\definecolor{accentlt}{HTML}{2E75B6}
\definecolor{rulegray}{HTML}{C3CEDA}
\definecolor{codebg}{HTML}{F6F8FA}
\definecolor{codeframe}{HTML}{D9E1EA}
\definecolor{cmt}{HTML}{6A737D}
\definecolor{kw}{HTML}{0B5394}
\definecolor{str}{HTML}{A31515}
\definecolor{num}{HTML}{098658}
\definecolor{noteblue}{HTML}{EAF2FB}
\definecolor{qabg}{HTML}{F3EEFA}
\definecolor{qarule}{HTML}{7E57C2}
\definecolor{warnbg}{HTML}{FBF1E6}
\definecolor{warnrule}{HTML}{D9822B}
\definecolor{okbg}{HTML}{E8F5E9}
\definecolor{okrule}{HTML}{2E7D32}

% ---- Tiêu đề mục ------------------------------------------------------------
\usepackage{titlesec}
\titleformat{\section}{\Large\bfseries\color{accent}}{\thesection}{0.6em}{}
\titleformat{\subsection}{\large\bfseries\color{accentlt}}{\thesubsection}{0.6em}{}
\titleformat{\subsubsection}{\normalsize\bfseries\color{accentlt}}{\thesubsubsection}{0.5em}{}
\titlespacing*{\section}{0pt}{1.4em}{0.6em}

% ---- Danh sách --------------------------------------------------------------
\usepackage{enumitem}
\setlist[itemize]{leftmargin=1.4em,topsep=2pt,itemsep=2pt}
\setlist[enumerate]{leftmargin=1.6em,topsep=2pt,itemsep=2pt}

% ---- Bảng -------------------------------------------------------------------
\usepackage{booktabs}
\usepackage{tabularx}
\usepackage{array}
\usepackage{colortbl}
\newcolumntype{L}[1]{>{\raggedright\arraybackslash}p{#1}}

% ---- Hộp callout (tcolorbox) ------------------------------------------------
\usepackage[most]{tcolorbox}
\newtcolorbox{notebox}{colback=noteblue,colframe=accentlt,boxrule=0pt,
  leftrule=3pt,arc=2pt,left=8pt,right=8pt,top=5pt,bottom=5pt}
\newtcolorbox{qabox}{colback=qabg,colframe=qarule,boxrule=0pt,
  leftrule=3pt,arc=2pt,left=8pt,right=8pt,top=5pt,bottom=5pt}
\newtcolorbox{warnbox}{colback=warnbg,colframe=warnrule,boxrule=0pt,
  leftrule=3pt,arc=2pt,left=8pt,right=8pt,top=5pt,bottom=5pt}

% ---- Mã nguồn (fvextra Verbatim — render đúng tiếng Việt dưới XeLaTeX) -------
% Lưu ý: listings reorder/garble dấu tiếng Việt trong code, nên dùng fvextra.
\usepackage{fvextra}
\fvset{
  numbers=left,
  numbersep=8pt,
  frame=single,
  framesep=6pt,
  rulecolor=\color{codeframe},
  fontsize=\footnotesize,
  breaklines=true,
  breakanywhere=true,
  breaksymbolleft={\footnotesize\color{cmt}\ensuremath{\hookrightarrow}},
}
% \code: định danh inline; underscore tự cho phép ngắt dòng (tránh tràn bảng).
\newcommand{\code}[1]{{\ttfamily\small\def\_{\textunderscore\discretionary{}{}{}}#1}}

% ---- Liên kết ---------------------------------------------------------------
\usepackage{hyperref}
\hypersetup{
  colorlinks=true, linkcolor=accent, urlcolor=accentlt, citecolor=accent,
  pdftitle={Phan cong code van dap - Nguyen Trong Phuc - X509Certificate},
  pdfauthor={Nguyen Trong Phuc},
}

% ---- Header / footer --------------------------------------------------------
\usepackage{fancyhdr}
\pagestyle{fancy}\fancyhf{}
\renewcommand{\headrulewidth}{0.4pt}
\renewcommand{\footrulewidth}{0pt}
\fancyhead[L]{\footnotesize\color{accent}X509Certificate — Phân công vấn đáp}
\fancyhead[R]{\footnotesize\color{accent}Nguyễn Trọng Phúc}
\fancyfoot[C]{\thepage}

% ---- Hộp nhỏ dùng cho sơ đồ (thay cho TikZ, biên dịch ổn định) --------------
\newcommand{\flowbox}[2]{\colorbox{#1}{\parbox{8.4cm}{\centering #2}}}
\newcommand{\stepbox}[1]{\fcolorbox{accent}{codebg}{\footnotesize\strut #1}}

\begin{document}

%==============================================================================
% TRANG BÌA
%==============================================================================
\begin{titlepage}
\centering
\vspace*{2.2cm}
{\color{accent}\rule{\linewidth}{1.2pt}}\\[0.6cm]
{\Huge\bfseries Phân công code vấn đáp đồ án}\\[0.35cm]
{\Huge\bfseries\color{accent} X509Certificate}\\[0.5cm]
{\Large Phần việc chi tiết của thành viên}\\[0.3cm]
{\LARGE\bfseries Nguyễn Trọng Phúc}\\[0.45cm]
{\large\itshape Lõi mật mã, Root CA\\ và Phát hành chứng chỉ}\\[0.5cm]
{\color{accent}\rule{\linewidth}{1.2pt}}\\[1.4cm]

\begin{minipage}{0.85\linewidth}
\centering\small
Tài liệu này mô tả chi tiết phần code, lý thuyết X.509 cần nắm, các test case
và cách tích hợp cho phần việc \textbf{lõi mật mã + Root CA + phát hành cert},
gắn trực tiếp với mã nguồn thực tế của đồ án trong các thư mục
\code{src/core/} (\code{ca.py}, \code{cert\_builder.py}, \code{csr.py}),
\code{src/services/} (\code{ca\_admin.py}, \code{csr\_admin.py},
\code{csr\_workflow.py}, \code{customer\_keys.py}) và \code{src/ui/}.
\end{minipage}

\vfill
{\large Môn: Mã hoá và Ứng dụng (MHUD)}\\[0.2cm]
{\large Học phần đồ án X.509 Certificate}\\[0.6cm]
{\small Ngày biên soạn: 14/06/2026}
\end{titlepage}

\tableofcontents
\thispagestyle{fancy}
\clearpage

%==============================================================================
\section{Mục tiêu và phạm vi phần việc}
%==============================================================================

Theo bảng phân công chung, thành viên \textbf{Nguyễn Trọng Phúc} phụ trách
\textbf{lõi mật mã, Root CA và phát hành chứng chỉ}. Đây là phần \emph{đứng đầu}
trong vòng đời chứng chỉ: nó tạo ra gốc tin cậy (Root CA self-signed), sinh cặp
khóa cho khách hàng, đóng gói yêu cầu cấp chứng chỉ (CSR theo chuẩn PKCS\#10),
và cuối cùng \textbf{phát hành} chứng chỉ X.509 v3 hợp lệ bằng cách dùng private
key của Root CA để ký. Nói cách khác, cụm này \emph{sản xuất} ra mọi vật phẩm
mật mã mà các cụm khác (kiểm tra chuỗi, thu hồi CRL \& OCSP) sẽ tiêu thụ về sau.

Phạm vi gồm bốn trục chính:
\begin{enumerate}
  \item \textbf{Sinh khóa và chữ ký số} — RSA 2048/3072/4096, ký SHA-256 với
        padding PKCS\#1 v1.5.
  \item \textbf{Root CA} — tạo cert self-signed, mã hóa private key
        \emph{encrypt-at-rest} (AES-256-GCM) trước khi lưu DB, rotate, publish
        ra trust store.
  \item \textbf{CSR (PKCS\#10)} — build/parse/verify; chứng minh quyền sở hữu
        khóa (proof of possession).
  \item \textbf{Phát hành cert} — duyệt CSR, build cert end-entity v3 với đầy đủ
        extension (KeyUsage, EKU, SAN, CRLDP, AIA, SKI/AKI), ký bằng Root CA.
\end{enumerate}

\subsection{Ánh xạ khái niệm $\to$ file code $\to$ vai trò}

\begin{table}[h]
\centering\small
\renewcommand{\arraystretch}{1.35}
\begin{tabularx}{\linewidth}{L{4.4cm} L{4.4cm} X}
\toprule
\textbf{Khái niệm} & \textbf{File / hàm code} & \textbf{Vai trò} \\
\midrule
Sinh keypair RSA + tạo cert & \code{core/cert\_builder.py} & Hàm dựng cert X.509 v3, builder pattern của thư viện \code{cryptography}. \\
Root CA self-signed (production) & \code{services/ca\_admin.py} \newline \code{\_generate\_root\_ca()} & Sinh Root CA, mã hóa key AES-GCM, lưu DB, rotate. \\
Root CA self-signed (legacy/lab) & \code{core/ca.py} \newline \code{load\_or\_create\_issuer()} & Bản cũ lưu cert/key ra file + trust store. \\
Sinh keypair khách hàng & \code{services/customer\_keys.py} \newline \code{generate\_keypair()} & RSA + encrypt-at-rest, AAD bind theo id, BOLA guard. \\
CSR (PKCS\#10) & \code{core/csr.py} & \code{build\_csr} / \code{parse\_csr} / \code{verify\_csr\_signature}. \\
Customer nộp CSR & \code{services/csr\_workflow.py} \newline \code{submit\_csr()} & Validate domain/wildcard, ký CSR, lưu pending. \\
Duyệt và phát hành cert & \code{services/csr\_admin.py} \newline \code{approve\_csr()} & Verify CSR $\to$ load Root CA $\to$ build cert $\to$ ký. \\
Tham số cấp phát & \code{services/system\_config.py} & key\_size / validity / algo / URL nhúng cert. \\
UI sinh/rotate Root CA & \code{ui/admin/root\_ca\_view.py} & Màn hình quản lý Root CA. \\
UI duyệt CSR $\to$ cấp cert & \code{ui/admin/csr\_queue\_view.py} & Queue duyệt, approve/reject. \\
UI khách hàng & \code{ui/customer/my\_keys\_view.py} \newline \code{ui/customer/csr\_submit\_view.py} & Sinh keypair, gửi CSR. \\
\bottomrule
\end{tabularx}
\end{table}

\begin{notebox}
\textbf{Điểm cần nhớ khi mở bài:} cụm này là \emph{bên phát hành} (issuer side).
Mọi thứ nó tạo ra đều có một chữ ký số làm gốc tin cậy: CSR được ký bằng khóa
khách hàng (chứng minh sở hữu), cert được ký bằng khóa Root CA (chứng minh do CA
phát hành), Root CA tự ký chính mình (vì nó là gốc, không có ai trên nó).
\end{notebox}

%==============================================================================
\section{Kiến trúc tổng thể và vị trí của thành viên trong hệ thống}
%==============================================================================

Hệ thống dùng mô hình \textbf{Root CA + Trust Store} hai tầng: một Root CA
self-signed ký trực tiếp các cert end-entity (không có intermediate CA), và
client tin cậy cert nhờ có Root CA trong trust store. Cụm của thành viên nằm ở
\emph{nửa trên} của vòng đời chứng chỉ.

\subsection{Luồng phát hành chứng chỉ (sơ đồ)}

\begin{center}
\stepbox{(1) Admin: tạo Root CA self-signed \quad ca\_admin.create\_root\_ca} \\[3pt]
$\downarrow$ \\[3pt]
\stepbox{(2) Customer: sinh keypair RSA \quad customer\_keys.generate\_keypair} \\[3pt]
$\downarrow$ \\[3pt]
\stepbox{(3) Customer: tạo + ký CSR (PKCS\#10) \quad csr.build\_csr / csr\_workflow.submit\_csr} \\[3pt]
$\downarrow$ \\[3pt]
\stepbox{(4) Admin: verify CSR $\to$ load Root CA $\to$ ký cert \quad csr\_admin.approve\_csr} \\[3pt]
$\downarrow$ \\[3pt]
\stepbox{(5) Cert X.509 v3 phát hành \quad cert\_builder.issue\_cert\_from\_csr}
\end{center}

\subsection{Phân lớp module và quan hệ}

\begin{table}[h]
\centering\small
\renewcommand{\arraystretch}{1.3}
\begin{tabularx}{\linewidth}{L{2.6cm} L{4.6cm} X}
\toprule
\textbf{Tầng} & \textbf{Module} & \textbf{Trách nhiệm} \\
\midrule
UI & \code{root\_ca\_view}, \code{csr\_queue\_view}, \code{my\_keys\_view}, \code{csr\_submit\_view} & Thu thập input, gọi service, hiển thị kết quả, ghi audit. \\
Service & \code{ca\_admin}, \code{csr\_admin}, \code{csr\_workflow}, \code{customer\_keys}, \code{system\_config} & Nghiệp vụ + giao dịch DB (atomic, BOLA guard, race-check). \\
Core (mật mã) & \code{cert\_builder}, \code{csr}, \code{ca}, \code{encryption} & Thao tác mật mã thuần (sinh khóa, ký, build cert, AES-GCM). \\
DB & \code{root\_ca}, \code{customer\_keys}, \code{csr\_requests}, \code{issued\_certs}, \code{system\_config} & Lưu trữ bền vững. \\
\bottomrule
\end{tabularx}
\end{table}

\begin{notebox}
\textbf{Quy tắc thiết kế:} tầng \code{core/} chỉ làm mật mã \emph{thuần}, không
chạm DB; tầng \code{services/} ghép mật mã với DB trong transaction; tầng UI
không bao giờ tự ký/giải mã mà luôn gọi xuống service. Nhờ vậy mỗi tầng test
được độc lập.
\end{notebox}

%==============================================================================
\section{Phần code cần làm (chi tiết)}
%==============================================================================

\subsection{Tạo Root CA self-signed (\code{ca\_admin.\_generate\_root\_ca})}

Root CA là điểm neo tin cậy. Vì nó là gốc nên \emph{tự ký} chính mình:
\code{subject == issuer}, và chữ ký được tạo bằng \emph{chính private key của
nó}. Đây là điểm khác biệt cốt lõi so với cert end-entity (do Root CA ký).

\begin{Verbatim}
def _generate_root_ca(common_name, key_size, validity_days):
    key = rsa.generate_private_key(public_exponent=65537, key_size=key_size)
    # subject == issuer  → đặc trưng của cert self-signed
    subject = issuer = x509.Name([
        x509.NameAttribute(NameOID.COUNTRY_NAME,      "VN"),
        x509.NameAttribute(NameOID.ORGANIZATION_NAME, "X509 Demo"),
        x509.NameAttribute(NameOID.COMMON_NAME,       common_name),
    ])
    now = datetime.now(timezone.utc)
    cert = (
        x509.CertificateBuilder()
        .subject_name(subject).issuer_name(issuer)
        .public_key(key.public_key())
        .serial_number(x509.random_serial_number())
        .not_valid_before(now - timedelta(minutes=1))
        .not_valid_after(now + timedelta(days=validity_days))
        # ... các extension bên dưới ...
        .sign(key, hashes.SHA256())   # ← TỰ KÝ bằng chính khóa của nó
    )
    return cert, key
\end{Verbatim}

Hai extension định danh \enquote{đây là CA}:

\begin{Verbatim}
.add_extension(
    x509.BasicConstraints(ca=True, path_length=0), critical=True,
)
.add_extension(
    x509.KeyUsage(
        digital_signature=True, key_cert_sign=True, crl_sign=True,
        key_encipherment=False, content_commitment=False,
        data_encipherment=False, key_agreement=False,
        encipher_only=False, decipher_only=False,
    ),
    critical=True,
)
.add_extension(
    x509.SubjectKeyIdentifier.from_public_key(key.public_key()),
    critical=False,
)
\end{Verbatim}

Giải thích từng ý:
\begin{itemize}
  \item \code{BasicConstraints(ca=True, path\_length=0)}: xác nhận đây là CA và
        \code{path\_length=0} giới hạn chuỗi — Root CA chỉ được ký trực tiếp
        end-entity cert, \emph{không} được ký thêm CA con (intermediate). Đặt
        \code{critical=True} để client buộc phải hiểu và tuân thủ.
  \item \code{KeyUsage}: bật \code{key\_cert\_sign} (được ký cert) và
        \code{crl\_sign} (được ký CRL) — đúng hai vai trò của một CA. Tắt
        \code{key\_encipherment} vì khóa CA không dùng để mã hóa dữ liệu.
  \item \code{SubjectKeyIdentifier} (SKI): vân tay của public key Root CA; cert
        con sẽ trỏ về SKI này qua \code{AuthorityKeyIdentifier} (AKI), giúp
        client nối chuỗi nhanh và chính xác.
\end{itemize}

\subsection{Lưu Root CA encrypt-at-rest (\code{ca\_admin.create\_root\_ca})}

Private key của Root CA là tài sản nhạy cảm nhất hệ thống — lộ nó là lộ toàn bộ
tin cậy. Vì vậy nó được \textbf{mã hóa AES-256-GCM} trước khi ghi DB, và việc
ghi diễn ra atomic: deactivate Root CA cũ rồi insert cái mới trong cùng một
transaction.

\begin{Verbatim}
cert, key = _generate_root_ca(common_name, key_size, validity_days)
cert_pem = cert.public_bytes(serialization.Encoding.PEM)
key_pem  = _serialize_private_key_pem(key)   # PKCS8 PEM, không password

# Encrypt-at-rest. AAD = b"root_ca" để bind ciphertext với loại record.
nonce, ct = encrypt_blob(key_pem, aad=ROOT_CA_AAD)
# ...
with transaction(db_path) as conn:
    conn.execute("UPDATE root_ca SET is_active = 0 WHERE is_active = 1")
    cur = conn.execute(
        "INSERT INTO root_ca (common_name, serial_hex, cert_pem, "
        " encrypted_private_key, gcm_nonce, ..., is_active) "
        "VALUES (?, ?, ?, ?, ?, ..., 1)",
        (common_name, serial_hex, cert_pem, ct, nonce, ...),
    )
    new_id = cur.lastrowid
\end{Verbatim}

\begin{itemize}
  \item \textbf{Tại sao mã hóa:} kể cả khi file DB bị copy ra ngoài, kẻ tấn công
        cũng chỉ thấy ciphertext; không có master key thì không giải mã được
        private key.
  \item \textbf{AAD (Additional Authenticated Data) =} \code{b"root\_ca"}: dữ
        liệu bổ trợ được \emph{xác thực} (không mã hóa) bởi GCM. Nếu ai đó lấy
        ciphertext của Root CA rồi cố giải mã trong ngữ cảnh khác (AAD khác),
        tag check sẽ thất bại. Nó \enquote{cột} ciphertext vào đúng loại bản ghi.
  \item \textbf{Atomic:} \code{UPDATE ... is\_active=0} và \code{INSERT ...
        is\_active=1} cùng một transaction $\to$ không bao giờ tồn tại trạng thái
        không có hoặc có hai Root CA active.
\end{itemize}

\begin{warnbox}
Khi \textbf{rotate} (sinh Root CA mới đè cái cũ), các cert đã phát hành vẫn mang
chữ ký của Root CA \emph{cũ}. Client sẽ FAIL verify cho đến khi admin re-issue
cert mới. UI có cảnh báo đỏ về điều này (xem \code{root\_ca\_view}).
\end{warnbox}

\subsection{Nạp Root CA để ký (\code{load\_active\_root\_ca\_with\_key})}

Đây là hàm \emph{duy nhất} giải mã private key Root CA, chỉ gọi khi cần ký cert
mới hoặc ký CRL.

\begin{Verbatim}
def load_active_root_ca_with_key(db_path):
    with conn_scope(db_path) as conn:
        row = conn.execute(
            "SELECT cert_pem, encrypted_private_key, gcm_nonce "
            "FROM root_ca WHERE is_active = 1 LIMIT 1"
        ).fetchone()
    if row is None:
        raise CAError("Chưa có Root CA active. Admin cần tạo Root CA trước.")
    cert = x509.load_pem_x509_certificate(row["cert_pem"])
    key_pem = decrypt_blob(
        row["gcm_nonce"], row["encrypted_private_key"], aad=ROOT_CA_AAD,
    )
    key = serialization.load_pem_private_key(key_pem, password=None)
    return cert, key
\end{Verbatim}

Lưu ý AAD truyền vào \code{decrypt\_blob} phải đúng bằng AAD lúc encrypt
(\code{ROOT\_CA\_AAD}), nếu không GCM tag check fail và raise lỗi — cơ chế chống
nhầm/giả mạo ngữ cảnh.

\subsection{Bản legacy core/ca.py (file \& trust store)}

\code{core/ca.py} giữ đường legacy/lab: lưu cert/key Root CA ra \emph{file} thay
vì DB, và publish ra trust store cho client. Hàm \code{load\_or\_create\_issuer}
\enquote{load nếu có, sinh nếu chưa}. Vì là đường lab nên private key \emph{không}
encrypt-at-rest; bù lại nó siết quyền file ở mức OS.

\begin{Verbatim}
def load_or_create_issuer(cert_path, key_path):
    if os.path.exists(cert_path) and os.path.exists(key_path):
        with open(key_path, "rb") as f:
            key = serialization.load_pem_private_key(f.read(), password=None)
        with open(cert_path, "rb") as f:
            cert = x509.load_pem_x509_certificate(f.read())
        return cert, key
    # chưa có → sinh Root CA self-signed mới (subject == issuer)
    key = rsa.generate_private_key(public_exponent=65537, key_size=2048)
    # ... build cert + ghi file ...
    _restrict_key_file_permissions(key_path)  # chmod 0600 / icacls
    return cert, key

def publish_root_ca_to_trust_store(root_ca_cert, trust_store_dir):
    out_path = os.path.join(trust_store_dir, "root_ca.crt")
    with open(out_path, "wb") as f:
        f.write(root_ca_cert.public_bytes(serialization.Encoding.PEM))
    return out_path
\end{Verbatim}

\begin{notebox}
Phân biệt rõ khi vấn đáp: \code{ca\_admin.py} là đường \textbf{production} (DB +
AES-GCM); \code{core/ca.py} là đường \textbf{legacy/lab} (file + trust store).
Cả hai đều sinh Root CA self-signed với cùng bộ extension.
\end{notebox}

\subsection{Sinh keypair khách hàng (\code{customer\_keys.generate\_keypair})}

Mỗi khách hàng có nhiều keypair RSA. Private key cũng \emph{encrypt-at-rest},
nhưng AAD ở đây phụ thuộc \emph{id} của bản ghi — mà id chỉ biết sau khi INSERT.
Giải pháp: insert placeholder để lấy id, rồi encrypt với AAD đúng và UPDATE lại,
tất cả trong một transaction.

\begin{Verbatim}
def _aad_for(key_id):
    return f"customer_keys:{key_id}".encode("ascii")

def generate_keypair(owner_id, name, key_size, db_path):
    # ... validate name + key_size ...
    rsa_key  = rsa.generate_private_key(public_exponent=65537, key_size=key_size)
    priv_pem = _serialize_private_pem(rsa_key)   # PKCS8 PEM
    pub_pem  = _serialize_public_pem(rsa_key)
    with transaction(db_path) as conn:
        # 1) INSERT placeholder (ct/nonce rỗng) để lấy id
        cur = conn.execute(
            "INSERT INTO customer_keys (owner_id, name, algorithm, key_size, "
            " public_key_pem, encrypted_private_key, gcm_nonce, created_at) "
            "VALUES (?, ?, ?, ?, ?, ?, ?, ?)",
            (owner_id, name, "RSA", key_size, pub_pem, b"", b"", now),
        )
        key_id = cur.lastrowid
        # 2) Encrypt với AAD bind theo id  3) UPDATE ciphertext thật
        nonce, ct = encrypt_blob(priv_pem, aad=_aad_for(key_id))
        conn.execute(
            "UPDATE customer_keys SET encrypted_private_key = ?, gcm_nonce = ? "
            "WHERE id = ?", (ct, nonce, key_id),
        )
\end{Verbatim}

\textbf{BOLA guard:} mọi hàm đọc khóa đều thêm \code{WHERE id=? AND owner\_id=?}.
Service không tin caller có đúng quyền — luôn lọc theo chủ sở hữu để chống tham
chiếu đối tượng chéo (Broken Object Level Authorization / IDOR).

\begin{Verbatim}
def load_private_key(key_id, owner_id, db_path):
    row = conn.execute(
        "SELECT encrypted_private_key, gcm_nonce FROM customer_keys "
        "WHERE id = ? AND owner_id = ?", (key_id, owner_id),
    ).fetchone()
    if row is None:
        raise CustomerKeyError(f"Không tìm thấy keypair id={key_id} thuộc về bạn.")
    pem = decrypt_blob(row["gcm_nonce"], row["encrypted_private_key"],
                       aad=_aad_for(key_id))
    return serialization.load_pem_private_key(pem, password=None)
\end{Verbatim}

\subsection{Tạo CSR theo PKCS\#10 (\code{core/csr.build\_csr})}

CSR (Certificate Signing Request) là gói \enquote{xin cấp chứng chỉ}: chứa
subject, public key và SAN, được \emph{ký bằng private key chủ nhân}. Chữ ký này
là bằng chứng \textbf{proof of possession} — chứng minh người xin thực sự nắm
private key tương ứng với public key trong yêu cầu.

\begin{Verbatim}
def build_csr(private_key, common_name, san_list=None,
              organization="X509 Demo", country="VN"):
    common_name = common_name.strip()
    san_values = list(san_list or [])
    if common_name not in san_values:
        san_values.insert(0, common_name)  # TLS hiện đại: CN phải nằm trong SAN
    subject = x509.Name([
        x509.NameAttribute(NameOID.COUNTRY_NAME,      country),
        x509.NameAttribute(NameOID.ORGANIZATION_NAME, organization),
        x509.NameAttribute(NameOID.COMMON_NAME,       common_name),
    ])
    builder = (
        x509.CertificateSigningRequestBuilder()
        .subject_name(subject)
        .add_extension(
            x509.SubjectAlternativeName(_build_san_list(san_values)),
            critical=False,
        )
    )
    return builder.sign(private_key, hashes.SHA256())  # ← ký = proof of possession
\end{Verbatim}

Phía duyệt CSR phải verify chữ ký trước khi cấp cert. Hàm verify dùng public key
\emph{nằm trong chính CSR} để kiểm tra:

\begin{Verbatim}
def verify_csr_signature(csr):
    if hasattr(csr, "is_signature_valid"):     # cryptography >= 40
        return bool(csr.is_signature_valid)
    try:                                        # fallback: verify thủ công
        csr.public_key().verify(
            csr.signature, csr.tbs_certrequest_bytes, *_verify_padding_args(csr),
        )
        return True
    except InvalidSignature:
        return False

def _verify_padding_args(csr):
    # RSA cần padding PKCS1v15 + hash; EC chỉ cần hash
    if isinstance(csr.public_key(), rsa.RSAPublicKey):
        return (padding.PKCS1v15(), csr.signature_hash_algorithm)
    return (csr.signature_hash_algorithm,)
\end{Verbatim}

\subsection{Customer nộp CSR (\code{csr\_workflow.submit\_csr})}

Service customer ghép \code{customer\_keys} với \code{core/csr}: kiểm tra quyền
sở hữu key, validate tên miền (cho phép wildcard \code{*.}), tạo + ký CSR, lưu
DB với \code{status='pending'}.

\begin{Verbatim}
def _validate_common_name(cn):
    cn = (cn or "").strip()
    if not cn:               raise CSRError("Common Name không được rỗng.")
    if len(cn) > 253:        raise CSRError("Common Name dài quá 253 ký tự.")
    name_to_check = cn[2:] if cn.startswith("*.") else cn   # wildcard '*.'
    if not all(c.isalnum() or c in ".-" for c in name_to_check):
        raise CSRError("Common Name chỉ chấp nhận chữ/số/dấu '.' '-' ...")
    return cn

def submit_csr(requester_id, customer_key_id, common_name, san_list, db_path):
    common_name = _validate_common_name(common_name)
    meta = get_key_meta(customer_key_id, requester_id, db_path)  # BOLA guard
    if meta is None:
        raise CSRError(f"Keypair id={customer_key_id} không thuộc về bạn ...")
    key = load_private_key(customer_key_id, requester_id, db_path)
    csr = build_csr(key, common_name=common_name, san_list=san_list)
    csr_pem = csr_to_pem(csr)
    with transaction(db_path) as conn:
        cur = conn.execute(
            "INSERT INTO csr_requests (requester_id, customer_key_id, "
            " common_name, san_list_json, csr_pem, status, submitted_at) "
            "VALUES (?, ?, ?, ?, ?, 'pending', ?)", (...),
        )
\end{Verbatim}

\subsection{Phát hành cert (\code{csr\_admin.approve\_csr})}

Đây là transaction trung tâm của cụm. Trình tự: (1) đọc CSR và kiểm tra
\code{status=='pending'}; (2) parse + verify chữ ký CSR; (3) load Root CA active
(giải mã key); (4) build + ký cert; (5) INSERT \code{issued\_certs} và UPDATE
\code{csr\_requests} atomic, kèm re-check chống race.

\begin{Verbatim}
# Bước 2: parse + verify CSR (proof of possession)
csr_pem = bytes(row["csr_pem"])
csr_obj = parse_csr(csr_pem)
if not verify_csr_signature(csr_obj):
    raise CSRAdminError("Chữ ký CSR không hợp lệ — có thể bị sửa hoặc khớp "
                        "sai keypair. Từ chối phát hành.")

# Bước 3: load Root CA active (decrypt private key)
ca_cert, ca_key = load_active_root_ca_with_key(db_path)

# Bước 4: build + sign cert end-entity từ CSR
cert, serial_number = issue_cert_from_csr(
    csr_obj, ca_cert, ca_key, validity_days=validity_days,
    ocsp_url=ocsp_url, crl_url=crl_url,
)
cert_pem = cert.public_bytes(serialization.Encoding.PEM)
\end{Verbatim}

\begin{Verbatim}
# Bước 5: ghi DB atomic + re-check chống race
with transaction(db_path) as conn:
    cur_status = conn.execute(
        "SELECT status FROM csr_requests WHERE id = ?", (csr_id,),
    ).fetchone()
    if cur_status is None or cur_status["status"] != "pending":
        raise CSRAdminError("CSR đã được xử lý bởi admin khác. Hủy approve.")
    conn.execute(
        "INSERT INTO issued_certs (csr_request_id, owner_id, serial_hex, "
        " common_name, cert_pem, not_valid_before, not_valid_after, "
        " issued_at, issued_by) VALUES (?, ?, ?, ?, ?, ?, ?, ?, ?)", (...),
    )
    conn.execute(
        "UPDATE csr_requests SET status = 'approved', reviewed_at = ?, "
        " reviewed_by = ? WHERE id = ?", (now, admin_id, csr_id),
    )
\end{Verbatim}

\begin{itemize}
  \item \textbf{Verify trước, ký sau:} không bao giờ ký một CSR chưa được xác
        thực chữ ký — nếu không, kẻ tấn công có thể sửa subject của CSR người
        khác và xin cert giả.
  \item \textbf{Re-check trong transaction:} dù đã kiểm tra \code{pending} ở
        đầu, vẫn kiểm tra lại \emph{bên trong} transaction phòng trường hợp hai
        admin approve cùng lúc (race) — chỉ một người thắng.
  \item \textbf{Atomic INSERT+UPDATE:} cert và việc đổi trạng thái CSR phải cùng
        commit; không thể có cert được cấp mà CSR vẫn \code{pending}.
\end{itemize}

\subsection{Builder cert v3 (\code{cert\_builder.\_build\_end\_entity\_cert})}

Hàm internal dựng cert end-entity với bộ extension cố định của TLS server cert.
Nó được dùng chung cho \code{issue\_cert\_from\_csr} (cấp mới) và
\code{reissue\_cert\_for\_renewal} (gia hạn).

\begin{Verbatim}
builder = (
    x509.CertificateBuilder()
    .subject_name(subject_name)
    .issuer_name(ca_cert.subject)          # issuer = Root CA (KHÁC subject)
    .public_key(public_key)                # public key lấy từ CSR
    .serial_number(x509.random_serial_number())
    .not_valid_before(now - timedelta(minutes=1))
    .not_valid_after(now + timedelta(days=validity_days))
    .add_extension(x509.BasicConstraints(ca=False, path_length=None),
                   critical=True)          # end-entity: KHÔNG phải CA
    .add_extension(
        x509.KeyUsage(digital_signature=True, key_encipherment=True,
                      key_cert_sign=False, crl_sign=False, ...),
        critical=True,
    )
    .add_extension(
        x509.ExtendedKeyUsage([ExtendedKeyUsageOID.SERVER_AUTH]),
        critical=False,                    # EKU = serverAuth (TLS server)
    )
)
\end{Verbatim}

\begin{Verbatim}
builder = (
    builder
    .add_extension(x509.CRLDistributionPoints([                # CRLDP
        x509.DistributionPoint(
            full_name=[x509.UniformResourceIdentifier(crl_url)],
            relative_name=None, reasons=None, crl_issuer=None)]),
        critical=False)
    .add_extension(x509.AuthorityInformationAccess([           # AIA / OCSP
        x509.AccessDescription(
            access_method=AuthorityInformationAccessOID.OCSP,
            access_location=x509.UniformResourceIdentifier(ocsp_url))]),
        critical=False)
    .add_extension(                                             # SKI
        x509.SubjectKeyIdentifier.from_public_key(public_key), critical=False)
    .add_extension(                                             # AKI → Root CA
        x509.AuthorityKeyIdentifier.from_issuer_public_key(ca_cert.public_key()),
        critical=False)
)
cert = builder.sign(private_key=ca_private_key, algorithm=hashes.SHA256())
\end{Verbatim}

So sánh nhanh cấu hình cert CA và end-entity:

\begin{table}[h]
\centering\small
\renewcommand{\arraystretch}{1.3}
\begin{tabularx}{\linewidth}{L{4.2cm} X X}
\toprule
\textbf{Extension} & \textbf{Root CA} & \textbf{End-entity (server)} \\
\midrule
BasicConstraints & \code{ca=True, path\_length=0} & \code{ca=False} \\
KeyUsage & \code{key\_cert\_sign}, \code{crl\_sign} & \code{digital\_signature}, \code{key\_encipherment} \\
ExtendedKeyUsage & (không) & \code{serverAuth} \\
SAN & (không) & DNS/IP của domain \\
issuer vs subject & bằng nhau (self-signed) & issuer = Root CA \\
\bottomrule
\end{tabularx}
\end{table}

\code{issue\_cert\_from\_csr} chỉ là wrapper lấy subject, public key, SAN từ CSR
rồi gọi \code{\_build\_end\_entity\_cert}; còn \code{reissue\_cert\_for\_renewal}
giữ nguyên subject + public key của cert cũ (admin gia hạn mà không đụng private
key của khách hàng). Hàm \code{tamper\_cert\_pem} lật một bit trong vùng chữ ký
để demo cert bị sửa vẫn parse được nhưng verify thất bại.

\subsection{Tham số cấp phát (\code{system\_config})}

Các tham số mặc định (key size, validity, thuật toán, URL CRL \& OCSP) được lưu
dạng key-value và nhúng vào cert lúc phát hành. \code{set\_config} chỉ chấp nhận
key trong whitelist để chống typo/inject.

\begin{Verbatim}
DEFAULTS = {
    "sig_algorithm":         "RSA-SHA256",
    "hash_algorithm":        "SHA256",
    "default_key_size":      "2048",
    "default_validity_days": "365",
    "root_ca_validity_days": "3650",
    "prod_crl_url":          "http://localhost:8889/crl.pem",
    "prod_ocsp_url":         "http://localhost:8888/ocsp",
}
ALLOWED_KEYS = frozenset(DEFAULTS.keys())  # whitelist chống inject
\end{Verbatim}

\subsection{Toàn cảnh nhà máy sinh cert --- \code{cert\_builder.py}}

\begin{notebox}
\code{cert\_builder.py} là \textbf{nhà máy} sinh \emph{mọi} certificate trong hệ
thống. Phúc \emph{sở hữu} trọn file này; các cụm khác chỉ \emph{gọi vào} chứ
không tự dựng cert. Bảng dưới liệt kê đủ các lối tạo cert để không sót hàm nào,
kèm việc \emph{ai gọi} (làm rõ ranh giới với Kiệt-renew và Quân-Lab).
\end{notebox}

\begin{table}[h]
\centering\footnotesize
\renewcommand{\arraystretch}{1.3}
\begin{tabularx}{\linewidth}{L{5.0cm} X L{2.6cm}}
\toprule
\textbf{Hàm} & \textbf{Tạo ra gì} & \textbf{Ai gọi} \\
\midrule
\code{generate\_rsa\_keypair} & Cặp khóa RSA (chưa phải cert) & submit\_csr\_lan, Lab (Quân) \\
\code{issue\_cert\_from\_csr} $\to$ \code{\_build\_end\_entity\_cert} & Cert end-entity \emph{từ CSR} (luồng chính) & \code{approve\_csr} (Phúc) \\
\code{create\_server\_cert\_signed\_by\_ca} & Cert server do Root CA ký \emph{trực tiếp}, không qua CSR & Verification Lab (Quân) \\
\code{reissue\_cert\_for\_renewal} & Cấp lại cert giữ nguyên subject $+$ public key & \code{renew\_cert} (Kiệt) \\
\code{create\_self\_signed\_cert} & Cert tự ký (issuer $=$ subject) --- đường legacy/test & code legacy \\
\code{tamper\_cert\_pem} & Lật 1 bit chữ ký để demo cert giả mạo & Verification Lab (Quân) \\
\code{save\_cert\_and\_key} / \code{save\_cert} / \code{load\_cert} / \code{load\_private\_key} & I/O cert/key ra/từ file PEM & nhiều nơi \\
\bottomrule
\end{tabularx}
\caption{Mọi lối tạo cert đều nằm trong \code{cert\_builder.py} (Phúc sở hữu).}
\end{table}

\textbf{Điểm chung và khác nhau.} Tất cả hàm trên đều dùng cùng một
\code{x509.CertificateBuilder()} (builder pattern). Khác nhau ở hai chỗ:
\begin{itemize}
  \item \emph{Ai là issuer}: \code{create\_self\_signed\_cert} và Root CA dùng
        \code{issuer == subject} (tự ký); còn \code{create\_server\_cert\_signed\_by\_ca},
        \code{issue\_cert\_from\_csr}, \code{reissue\_cert\_for\_renewal} dùng
        \code{issuer = ca\_cert.subject} và ký bằng \emph{private key của Root CA}.
  \item \emph{Bộ extension}: cert server/end-entity có \code{BasicConstraints(ca=False)},
        \code{KeyUsage} (digitalSignature, keyEncipherment), \code{ExtendedKeyUsage}
        serverAuth, SAN, CRLDP, AIA, SKI/AKI; còn Root CA có \code{ca=True},
        \code{keyCertSign}, \code{crlSign}.
\end{itemize}

Trích hàm tạo server cert do Root CA ký trực tiếp (lối Verification Lab dùng):

\begin{Verbatim}
def create_server_cert_signed_by_ca(server_private_key, ca_cert, ca_private_key,
                                    common_name="localhost", dns_names=None,
                                    validity_days=365, expired=False):
    issuer = ca_cert.subject                       # issuer = Root CA (KHÔNG self-signed)
    # ... builder: subject = CN; BasicConstraints(ca=False);
    #     KeyUsage(digital_signature, key_encipherment); EKU serverAuth;
    #     SAN; CRLDistributionPoints; AIA/OCSP; SKI/AKI ...
    certificate = builder.sign(private_key=ca_private_key,   # ← Root CA ký
                               algorithm=hashes.SHA256())
    return certificate, serial_number
\end{Verbatim}

\begin{notebox}
\textbf{Vì sao có hai lối tạo server cert?} \code{issue\_cert\_from\_csr} dùng cho
luồng \emph{thật} (khách gửi CSR, admin duyệt --- public key đến \emph{từ CSR}).
\code{create\_server\_cert\_signed\_by\_ca} dùng cho \emph{Verification Lab}
(Quân) cần sinh nhanh nhiều cert mẫu với cả keypair do Lab tự tạo. Cả hai cùng
cho ra cert end-entity hợp lệ do Root CA ký --- chỉ khác nguồn gốc của public key.
\end{notebox}

\subsection{Lớp UI của cụm}

\begin{itemize}
  \item \code{root\_ca\_view.py}: hiển thị Root CA active + lịch sử; nút
        \enquote{Sinh Root CA} / \enquote{rotate} / \enquote{Publish ra Trust
        Store}. Modal có cảnh báo đỏ khi rotate.
  \item \code{csr\_queue\_view.py}: queue duyệt CSR; \code{ApproveCSRDialog} gọi
        \code{approve\_csr} và ghi hai audit (\code{CSR\_APPROVED} +
        \code{CERT\_ISSUED}); \code{RejectCSRDialog} bắt buộc nhập lý do.
  \item \code{system\_config\_view.py}: form 5 trường, validate int, ghi audit
        các thay đổi.
  \item \code{my\_keys\_view.py}: sinh/xem/xóa keypair (xóa bị từ chối nếu key
        đang được CSR tham chiếu).
  \item \code{csr\_submit\_view.py}: chọn keypair, nhập CN + SAN, submit CSR
        (offline qua DB hoặc remote qua LAN tới máy Admin).
\end{itemize}

%==============================================================================
\section{Cần nắm để vấn đáp}
%==============================================================================

\begin{qabox}
\textbf{Hỏi:} Chữ ký số RSA hoạt động thế nào trong đồ án? Vì sao dùng padding
PKCS\#1 v1.5?\\[2pt]
\textbf{Đáp:} Ký = băm dữ liệu bằng SHA-256, rồi dùng private key \enquote{mã
hóa} (toán RSA) lên giá trị băm đã đệm (padding). Verify = dùng public key giải
ra giá trị băm, băm lại dữ liệu và so sánh. Đồ án dùng padding \code{PKCS1v15}
(thấy trong \code{\_verify\_padding\_args}) vì đây là sơ đồ chuẩn, tương thích
rộng và đủ cho mục tiêu học thuật; PSS hiện đại hơn nhưng PKCS1v15 đơn giản, dễ
giải thích. EC thì không cần padding, chỉ cần hash.
\end{qabox}

\begin{qabox}
\textbf{Hỏi:} Root CA là self-signed — vì sao nó tự ký được và điều đó có an
toàn không?\\[2pt]
\textbf{Đáp:} Root CA là \emph{gốc} của cây tin cậy, không có ai \enquote{trên}
nó để ký, nên nó tự ký: \code{subject == issuer} và chữ ký tạo bằng \emph{chính}
private key của nó (\code{.sign(key, SHA256())}). Bản thân việc self-signed
không tạo tin cậy — tin cậy đến từ việc cert đó được đặt vào \textbf{trust
store} của client. Client tin Root CA vì \enquote{quyết định} tin, không phải vì
chữ ký. An toàn nằm ở chỗ private key Root CA được bảo vệ (encrypt-at-rest) và
trust store được phân phối an toàn.
\end{qabox}

\begin{qabox}
\textbf{Hỏi:} CSR theo PKCS\#10 chứa gì? Proof of possession là gì và chống được
tấn công nào?\\[2pt]
\textbf{Đáp:} CSR chứa subject (CN, O, C), public key, SAN, và một \emph{chữ
ký} được tạo bằng private key chủ nhân lên toàn bộ phần TBS. Verify chữ ký bằng
chính public key trong CSR chính là \textbf{proof of possession} — chứng minh
người gửi thực sự sở hữu private key. Nó chống tấn công \enquote{xin cert cho
khóa không phải của mình}: nếu kẻ xấu lấy public key người khác và cố xin cert,
hắn không thể tạo chữ ký hợp lệ vì không có private key tương ứng.
\end{qabox}

\begin{qabox}
\textbf{Hỏi:} Khi cấp cert, các extension BasicConstraints / KeyUsage / EKU đặt
thế nào và vì sao?\\[2pt]
\textbf{Đáp:} Cert end-entity đặt \code{BasicConstraints(ca=False)} (không được
ký cert khác), \code{KeyUsage} bật \code{digital\_signature} +
\code{key\_encipherment} (đúng vai trò TLS), tắt \code{key\_cert\_sign}; và
\code{ExtendedKeyUsage = serverAuth} để khai báo \enquote{dùng làm cert server
TLS}. Ngược lại Root CA đặt \code{ca=True, path\_length=0} và bật
\code{key\_cert\_sign} + \code{crl\_sign}. Đặt critical cho BasicConstraints và
KeyUsage để client buộc phải tuân thủ.
\end{qabox}

\begin{qabox}
\textbf{Hỏi:} Vì sao phải encrypt-at-rest private key của CA? AAD dùng để làm
gì?\\[2pt]
\textbf{Đáp:} Vì nếu file DB bị lộ, private key plaintext sẽ cho phép kẻ tấn
công ký cert giả mạo bất kỳ — phá vỡ toàn bộ tin cậy. Mã hóa AES-256-GCM đảm bảo
chỉ ai có master key mới giải được. AAD (Additional Authenticated Data) là dữ
liệu được \emph{xác thực nhưng không mã hóa}; với Root CA dùng \code{b"root\_ca"},
với customer key dùng \code{b"customer\_keys:\{id\}"}. AAD \enquote{cột}
ciphertext vào đúng ngữ cảnh: nếu lấy ciphertext giải mã ở ngữ cảnh khác (AAD
khác hoặc id khác), GCM tag check thất bại $\to$ chống đổi chỗ/nhầm bản ghi.
\end{qabox}

\begin{qabox}
\textbf{Hỏi:} Proof of possession chống được gì cụ thể trong luồng cấp cert?\\[2pt]
\textbf{Đáp:} Nó chống việc cấp cert cho một public key mà người xin không nắm
private key — tức ngăn \enquote{cướp danh tính khóa}. Trong \code{approve\_csr},
nếu \code{verify\_csr\_signature} trả về False (CSR bị sửa subject/SAN, hoặc chữ
ký không khớp public key trong CSR), admin từ chối phát hành ngay. Nhờ vậy cert
chỉ được cấp cho đúng chủ thực sự của cặp khóa.
\end{qabox}

\begin{qabox}
\textbf{Hỏi:} Chuỗi renew (gia hạn) cert diễn ra thế nào? Admin có cần private
key của khách hàng không?\\[2pt]
\textbf{Đáp:} Gia hạn dùng \code{reissue\_cert\_for\_renewal}: nó giữ
\emph{nguyên} subject và public key của cert cũ, chỉ kéo dài thời hạn và sinh
serial mới, rồi ký lại bằng Root CA \emph{đang active}. Admin \textbf{không} cần
và không có private key của khách hàng — chỉ cần public key (đã nằm trong cert
cũ). Đây là lý do tách \code{\_build\_end\_entity\_cert} dùng chung cho cả cấp
mới lẫn gia hạn.
\end{qabox}

\begin{qabox}
\textbf{Hỏi:} SKI và AKI để làm gì trong cert?\\[2pt]
\textbf{Đáp:} SubjectKeyIdentifier (SKI) là vân tay của public key chủ thể;
AuthorityKeyIdentifier (AKI) là vân tay của public key bên ký (Root CA). Khi
client dựng chuỗi, nó khớp AKI của cert con với SKI của cert cha để nối nhanh và
chính xác, kể cả khi có nhiều CA cùng tên.
\end{qabox}

%==============================================================================
\section{Test case và demo}
%==============================================================================

\subsection{Bảng test case chính}

\begin{table}[h]
\centering\small
\renewcommand{\arraystretch}{1.3}
\begin{tabularx}{\linewidth}{L{4.6cm} X L{3.0cm}}
\toprule
\textbf{Tình huống} & \textbf{Kỳ vọng} & \textbf{File test} \\
\midrule
Tạo Root CA lần đầu & \code{is\_active=1}, cert PEM parse được, subject chứa CN & \code{test\_m4\_ca\_admin} \\
Private key không lưu raw vào DB & DB chỉ chứa ciphertext + nonce; decrypt mới ra key & \code{test\_m4\_ca\_admin} \\
Rotate Root CA & Root CA cũ \code{is\_active=0}, cái mới \code{=1} & \code{test\_m4\_ca\_admin} \\
Approve CSR happy path & cert ký bởi Root CA, public key khớp CSR & \code{test\_m6\_csr\_admin} \\
Approve khi chưa có Root CA & raise lỗi & \code{test\_m6\_csr\_admin} \\
CSR bị sửa chữ ký & verify fail $\to$ từ chối cấp & \code{test\_m6\_csr\_admin} \\
Race 2 admin approve & 1 thắng, 1 fail & \code{test\_m6\_csr\_admin} \\
SAN từ CSR sang cert & cert giữ đúng SAN & \code{test\_m6\_csr\_admin} \\
Sinh keypair + load lại & decrypt đúng key; BOLA: owner khác load fail & \code{test\_m5\_customer} \\
\bottomrule
\end{tabularx}
\end{table}

\subsection{Đoạn unit test gợi ý}

\begin{Verbatim}
def test_approve_csr_happy_path():
    env = TestEnv(with_root_ca=True)        # đã tạo Root CA trong setup
    csr_id = env.make_csr("demo.com", san_list=["www.demo.com"])
    issued = approve_csr(csr_id, env.admin_id, validity_days=365,
                         db_path=env.db_path)
    detail = get_csr_detail(csr_id, env.db_path)
    assert detail["status"] == "approved"
    # issuer của cert = subject của Root CA active
    ca_cert, _ = load_active_root_ca_with_key(env.db_path)
    cert = x509.load_pem_x509_certificate(<cert_pem từ issued_certs>)
    assert cert.issuer == ca_cert.subject
    # chữ ký cert verify được bằng public key Root CA
    ca_cert.public_key().verify(
        cert.signature, cert.tbs_certificate_bytes,
        padding.PKCS1v15(), cert.signature_hash_algorithm,
    )

def test_tampered_csr_rejected():
    env = TestEnv(with_root_ca=True)
    csr_id = env.make_csr("demo.com")
    # giả lập sửa CSR PEM trong DB → verify_csr_signature phải fail
    # → approve_csr raise CSRAdminError
\end{Verbatim}

\subsection{Lệnh chạy test}

\begin{Verbatim}
# Chạy từng module test của cụm
python tests/test_m4_ca_admin.py     # Root CA admin
python tests/test_m6_csr_admin.py    # duyệt CSR + phát hành cert
python tests/test_m5_customer.py     # keypair + CSR phía customer
\end{Verbatim}

\subsection{Kịch bản demo end-to-end}

\begin{enumerate}
  \item Admin mở \code{root\_ca\_view} $\to$ \enquote{Sinh Root CA} (RSA-2048,
        3650 ngày) $\to$ Publish ra Trust Store.
  \item Customer mở \code{my\_keys\_view} $\to$ \enquote{Sinh keypair mới}
        (RSA-2048).
  \item Customer mở \code{csr\_submit\_view} $\to$ chọn keypair, nhập
        \code{example.com} + SAN $\to$ Submit CSR (status \code{pending}).
  \item Admin mở \code{csr\_queue\_view} $\to$ chọn CSR $\to$ \enquote{Xem chi
        tiết} (CSR PEM) $\to$ \enquote{Approve} (365 ngày) $\to$ cert được phát
        hành, status chuyển \code{approved}.
  \item (Tùy chọn) Dùng \code{tamper\_cert\_pem} để minh họa cert bị sửa $\to$
        cụm kiểm tra chuỗi sẽ FAIL verify.
\end{enumerate}

%==============================================================================
\section{Tích hợp vào chương trình chung}
%==============================================================================

Cụm này là \emph{nhà sản xuất} cho phần còn lại của hệ thống. Quan hệ gọi/ghép
như sau:

\begin{itemize}
  \item \textbf{UI $\to$ Service:} \code{root\_ca\_view} gọi \code{create\_root\_ca},
        \code{publish\_active\_to\_trust\_store}; \code{csr\_queue\_view} gọi
        \code{approve\_csr} / \code{reject\_csr}; \code{csr\_submit\_view} gọi
        \code{submit\_csr}; \code{my\_keys\_view} gọi \code{generate\_keypair}.
  \item \textbf{Service $\to$ Core:} \code{csr\_admin.approve\_csr} ghép ba mảnh:
        \code{core/csr.verify\_csr\_signature},
        \code{ca\_admin.load\_active\_root\_ca\_with\_key} và
        \code{cert\_builder.issue\_cert\_from\_csr}.
  \item \textbf{Service $\to$ Config:} validity/URL CRL \& OCSP lấy từ
        \code{system\_config} để nhúng vào cert; URL còn fallback về
        \code{infra\_manager} để tự khớp port runtime.
  \item \textbf{Đầu ra cho cụm khác:} cert trong \code{issued\_certs} và Root CA
        trong trust store là input cho cụm \textbf{kiểm tra chuỗi}
        (\code{core/verify.py}); serial của cert là input cho cụm \textbf{thu
        hồi CRL \& OCSP}.
  \item \textbf{Audit:} mọi hành động nhạy cảm (tạo/rotate Root CA, approve/reject
        CSR, cấp cert, sinh keypair, đổi config) đều ghi log qua
        \code{services.audit.write\_audit}.
\end{itemize}

\begin{notebox}
Sợi chỉ xuyên suốt: \textbf{một chữ ký số ở mỗi mắt xích}. CSR mang chữ ký khách
hàng (proof of possession) $\to$ cert mang chữ ký Root CA (chứng thực) $\to$ Root
CA tự ký (gốc tin cậy). Nắm sợi chỉ này là nắm cả cụm.
\end{notebox}

%==============================================================================
\section{Bổ sung sau rà soát: các hàm/feature dễ bị hỏi}

\subsection{Nộp CSR qua LAN \& khái niệm khóa ``public-only''}

\begin{warnbox}
Câu hỏi điển hình: \emph{``Nộp CSR từ máy khác thì private key có bị gửi đi không?''}
--- KHÔNG. Ở chế độ \textbf{remote (LAN)}: client \code{build\_csr} \emph{tại máy
mình}, chỉ gửi \emph{CSR PEM} $+$ username/password tới Admin API. Máy Admin lưu
thành keypair \textbf{``public-only''} (cột \code{encrypted\_private\_key} rỗng).
\end{warnbox}

\begin{itemize}
  \item Hai chế độ submit: \emph{offline} (\code{submit\_csr} ghi thẳng DB, key nằm
        trong DB cùng máy) vs \emph{remote LAN} (key ở lại client).
  \item Hệ quả cờ \textbf{\code{is\_public\_only}}: \code{list\_keys}/\code{get\_key\_meta}
        gắn cờ này; \code{load\_private\_key} sẽ \emph{raise} với key public-only
        (``private key nằm trên máy client gốc'') --- không thể ký bằng key không có.
\end{itemize}

\subsection{Các hàm CSR/CA/cấu hình còn lại}

\begin{itemize}
  \item \textbf{\code{cancel\_csr} (customer tự hủy).} Chỉ hủy được CSR \code{pending};
        thực chất set \code{status='rejected'}, \code{reason='cancelled by requester'}
        (giữ lịch sử thay vì xóa cứng), có guard \code{WHERE requester\_id=?}; bị
        chặn ở chế độ remote (CSR nằm trên máy Admin).
  \item \textbf{\code{list\_my\_csr} / \code{get\_my\_csr\_by\_id} (BOLA).} Customer chỉ
        thấy CSR của mình (lọc \code{requester\_id}) --- đối lập với
        \code{list\_all\_csr}/\code{get\_csr\_detail} của admin (xem MỌI CSR vì là
        quyền quản trị).
  \item \textbf{\code{get\_active\_root\_ca} vs \code{load\_active\_root\_ca\_with\_key}.}
        \code{get\_active\_root\_ca} chỉ đọc metadata$+$\code{cert\_pem},
        \emph{KHÔNG} giải mã private key (giảm bề mặt lộ khóa); chỉ
        \code{load\_active\_root\_ca\_with\_key} mới decrypt --- và chỉ khi cần
        \emph{ký}. \code{list\_root\_ca\_history} liệt kê CA active$+$retired cho bảng
        lịch sử rotate.
  \item \textbf{Đọc cấu hình.} \code{get\_config}/\code{get\_all\_config} đọc giá trị;
        \code{get\_int\_config} có \emph{fallback} an toàn khi thiếu/parse lỗi;
        \code{seed\_defaults} idempotent lúc init DB. UI lấy default key\_size/validity
        qua \code{get\_int\_config}; \code{csr\_admin.\_default\_crl\_url}/\code{\_default\_ocsp\_url}
        ưu tiên config rồi fallback \code{infra\_manager}.
  \item \textbf{Trích thông tin CSR \& trust store.} \code{get\_csr\_common\_name}
        (CN, trả '' nếu trống) và \code{get\_csr\_san\_dns} (list DNS-SAN, rỗng nếu
        thiếu) --- dùng hiển thị/đối chiếu khi duyệt. \code{ca.load\_trust\_store}
        đọc mọi \code{.crt}/\code{.pem} trong trust store (bỏ qua file PEM hỏng thay
        vì crash), thiết kế list để mở rộng nhiều Root CA --- là đầu vào cho cụm verify.
\end{itemize}

%==============================================================================
\section{Checklist chuẩn bị vấn đáp}
%==============================================================================

\begin{itemize}
  \item[\textbf{$\square$}] Giải thích được tại sao Root CA self-signed
        (\code{subject==issuer}, tự ký bằng khóa của nó) và tin cậy đến từ trust
        store.
  \item[\textbf{$\square$}] Đọc trôi bộ extension Root CA: \code{BasicConstraints
        (ca=True, path\_length=0)}, \code{KeyUsage(key\_cert\_sign, crl\_sign)},
        SKI.
  \item[\textbf{$\square$}] Phân biệt cert CA và end-entity qua bảng so sánh
        extension (BasicConstraints, KeyUsage, EKU, SAN).
  \item[\textbf{$\square$}] Trình bày proof of possession: CSR ký bằng khóa chủ,
        admin verify trước khi cấp.
  \item[\textbf{$\square$}] Giải thích encrypt-at-rest (AES-256-GCM) và vai trò
        của AAD (\code{root\_ca} / \code{customer\_keys:\{id\}}).
  \item[\textbf{$\square$}] Mô tả 5 bước của \code{approve\_csr} và lý do
        verify-trước-ký-sau + re-check chống race + atomic.
  \item[\textbf{$\square$}] Nói được BOLA guard: luôn \code{WHERE owner\_id=?}.
  \item[\textbf{$\square$}] Giải thích chữ ký RSA + padding PKCS\#1 v1.5 +
        SHA-256.
  \item[\textbf{$\square$}] Trình bày chuỗi renew dùng
        \code{reissue\_cert\_for\_renewal} (giữ public key, admin không cần
        private key khách hàng).
  \item[\textbf{$\square$}] Cảnh báo rotate Root CA làm cert cũ FAIL verify đến
        khi re-issue.
  \item[\textbf{$\square$}] Chạy được 3 file test (\code{m4}, \code{m6},
        \code{m5}) và demo end-to-end tạo Root CA $\to$ cấp cert.
\end{itemize}

\end{document}
